\documentclass[11pt]{article}

\usepackage[T1]{fontenc}
\usepackage{lmodern}
\usepackage{amsmath,amssymb,amsthm}

\newcommand{\bplus}{\boxplus}
\newcommand{\bigbplus}{\mathop{\boxplus}}
\DeclareMathOperator{\hdet}{hdet}

\newtheorem{theorem}{Theorem}

\begin{document}

\title{A Cubic Identity in Hyperfields}
\author{}
\date{}
\maketitle

\begin{abstract}
We prove that, for every commutative hyperfield $\mathbb H$, if
$0\in x\bplus y\bplus z$, then
\[
0\in x^3\bplus y^3\bplus z^3\bplus (-xyz)\bplus (-xyz)\bplus (-xyz).
\]
The proof is a direct application of the standard determinant criterion for
singular matrices over hyperfields.
\end{abstract}

\section*{Preliminaries}

Throughout, a hyperfield means a commutative hyperfield. Finite hyper-sums are
unambiguous because the additive structure of a hyperfield is a canonical
commutative hypergroup. For an $n\times n$ matrix $A=(a_{ij})$ over a hyperfield,
define its hyperdeterminant by
\[
\hdet(A)
 :=
 \bigbplus_{\sigma\in S_n}
 \operatorname{sgn}(\sigma)\,
 a_{1,\sigma(1)}a_{2,\sigma(2)}\cdots a_{n,\sigma(n)}.
\]
Here $\operatorname{sgn}(\sigma)\in\{1,-1\}\subset \mathbb H$, and multiplication
by $-1$ means additive negation.

We use the following standard theorem of Viro.

\begin{theorem}[Viro's singularity criterion]\label{thm:viro}
Let $A=(a_{ij})\in M_n(\mathbb H)$. Suppose there exists a nonzero vector
$v=(v_1,\dots,v_n)^t\in \mathbb H^n$ such that, for every $i$,
\[
0\in \bigbplus_{j=1}^n a_{ij}v_j.
\]
Then
\[
0\in \hdet(A).
\]
Equivalently, hyperlinear dependence of the columns of a square matrix implies
singularity, where singularity means that $0$ belongs to the hyperdeterminant.
\end{theorem}

\begin{proof}[Reference]
This is the ``linear dependence implies singular determinant'' direction of
Viro's determinant theorem for matrices over hyperfields; see O.~Viro,
\emph{Hyperfields for tropical geometry I. Hyperfields and dequantization},
arXiv:1006.3034, Section~4, where it is proved in the discussion of singular
matrices over hyperfields. The same statement is also standard in the language
of represented matroids over hyperfields; see M.~Baker and N.~Bowler,
\emph{Matroids over hyperfields}, Adv. Math. 343 (2019), 821--863.
\end{proof}

\section*{The cubic identity}

\begin{theorem}
Let $\mathbb H$ be a hyperfield and let $x,y,z\in \mathbb H$. If
\[
0\in x\bplus y\bplus z,
\]
then
\[
0\in x^3\bplus y^3\bplus z^3\bplus (-xyz)\bplus (-xyz)\bplus (-xyz).
\]
\end{theorem}

\begin{proof}
Consider the circulant matrix
\[
M=
\begin{pmatrix}
x & y & z\\
z & x & y\\
y & z & x
\end{pmatrix}.
\]
Let $e=(1,1,1)^t$. Since $0\in x\bplus y\bplus z$ and hyperaddition is
commutative, each coordinate of $Me$ contains $0$:
\[
0\in x\bplus y\bplus z,\qquad
0\in z\bplus x\bplus y,\qquad
0\in y\bplus z\bplus x.
\]
Thus $0\in Me$. Also $e\neq 0$ because $1\neq 0$ in a hyperfield. By
Theorem~\ref{thm:viro},
\[
0\in \hdet(M).
\]

It remains only to expand $\hdet(M)$. The even permutations in $S_3$ give
\[
x\cdot x\cdot x=x^3,\qquad
y\cdot y\cdot y=y^3,\qquad
z\cdot z\cdot z=z^3,
\]
while the three odd permutations each give the product $xyz$, with a minus sign.
Hence
\[
\hdet(M)
=
x^3\bplus y^3\bplus z^3
\bplus (-xyz)\bplus (-xyz)\bplus (-xyz).
\]
Therefore
\[
0\in x^3\bplus y^3\bplus z^3
\bplus (-xyz)\bplus (-xyz)\bplus (-xyz),
\]
as required.
\end{proof}

\begin{thebibliography}{9}

\bibitem{Viro}
O.~Viro,
\emph{Hyperfields for tropical geometry I. Hyperfields and dequantization},
arXiv:1006.3034.

\bibitem{BakerBowler}
M.~Baker and N.~Bowler,
\emph{Matroids over hyperfields},
Adv. Math. \textbf{343} (2019), 821--863.

\end{thebibliography}

\end{document}